
en notant que @@M0@@ est un morphisme \emph{parfait}, car déduit par le changement de base plat @@M1@@ du morphisme @@M2@@, qui est parfait grâce à l'hypothèse @@M3@@ régulier. On trouve ainsi un groupe quotient
@@M4@@
de @@M5@@. Ce dernier est d'ailleurs de façon naturelle un module sur @@M6@@, car on constate aussitôt sur les définitions que le produit d'un élément de @@M7@@ par un élément de @@M8@@ est algébriquement équivalent à zéro si l'un des deux facteurs l'est, en utilisant (IV 2.12.2) pour @@M9@@.

On définit (quand on travaille avec le @@M10@@ naïf) l'idéal de @@M11@@ des éléments algébriquement équivalents à zéro comme étant l'idéal engendré par les @@M12@@, où @@M13@@ et @@M14@@ est algébriquement équivalent à zéro; l'anneau gradué quotient
@@M15@@
n'est alors autre que l'anneau de Chern associé au @@M16@@-anneau @@M17@@.

Enfin on définit le sous-groupe gradué de @@M18@@ des éléments algébriquement équivalents à zéro (@@M19@@ étant supposé localement noethérien), comme celui dont la composante de dimension @@M20@@ est l'image dans @@M21@@ du sous-groupe de @@M22@@ formé des éléments algébriquement équivalents à zéro dans @@M23@@. On trouve aussitôt que le groupe gradué quotient
@@M24@@
est un module sur l'anneau @@M25@@, par passage au quotient à partir de la structure de module de @@M26@@ sur @@M27@@ (3.1).

\textbf{7.11.6.} Les définitions qui précèdent garderaient un sens si on ne suppose pas @@M28@@ algébriquement clos, en prenant ci-dessus @@M29@@ \emph{géométriquement} connexe sur @@M30@@. Il n'est pas clair cependant que ces définitions seraient bien utiles; notamment, il ne semble pas que dans le cas du groupe @@M31@@, on
